\clearpage
\fancyhead[L]{\small BODE. ASTRONOMISCHES JAHRBUCH.}
\hypertarget{native-635-lower}{}
\begin{center}\small Göttingische gelehrte Anzeigen. Stück 2. Seiten 17..24. 3. Januar 1824.\\[3pt]\rule{.72\linewidth}{.3pt}\end{center}

\emph{Astronomisches Jahrbuch für das Jahr \textup{1826} mit einer Sammlung der neuesten Abhandlungen, Beobachtungen und Nachrichten aus den astronomischen Wissenschaften. Berechnet und herausgegeben von} Dr. J. E. \textsc{Bode}, \emph{königlicher Astronom usw.} Berlin 1823. Erhältlich beim Verfasser und in Kommission bei Ferd. Dümmler. 256 Seiten im Oktavformat, dazu eine Kupfertafel.

Dieser Jahrgang eröffnet die zweiten fünfzig Bände einer Reihe, die viel zur Verbreitung astronomischer Kenntnisse in Deutschland beigetragen hat und der wir eine noch lange Fortsetzung \SourceLeaf{636} unter der Leitung ihres würdigen Herausgebers wünschen. Unter den Himmelsereignissen des Jahres 1826, die von allgemeinerem Interesse sind, heben wir eine totale Mondfinsternis (am 14. November), eine partielle Sonnenfinsternis (am 29. November) und eine Bedeckung des Saturn durch den Mond (am 16. Februar) hervor.

Die Abhandlungen und Nachrichten im Anhang nehmen 168 Seiten ein. Zuerst einige Ergebnisse aus der Triangulierung im Königreich Hannover von Herrn Hofrat \textsc{Gauss}. Hier werden zunächst die geografischen Positionen von 22 ausgewählten Punkten mitgeteilt, wie sie im Herbst 1822 aus dem vorläufigen Anschluss an die \textsc{von Zach}'sche Basis abgeleitet wurden. Wir können hier hinzufügen, dass der inzwischen ausgeführte Anschluss an die von Herrn \textsc{Schumacher} in Holstein gemessene Grundlinie diese Ergebnisse nicht merklich verändert. Lediglich die weiter nördlich liegenden Breiten verringern sich um eine Zehntelbogensekunde. Die drei Türme in Hildesheim, die in diesem Verzeichnis ohne Namen vorkommen, lassen sich jetzt benennen. Der erste ist der Andreas-, der zweite der Michaelis- und der dritte der Jacobi-Turm. Dann folgt die Übersicht über die terrestrischen Refraktionen, die aus den gegenseitigen Zenitdistanzen entlang von 28 Linien des Dreiecksnetzes abgeleitet wurden. Das mittlere Verhältnis der Krümmung des terrestrischen Bogens zur beobachteten gesamten Refraktion beträgt 1 zu 0.1306. Durch die im Jahr 1823 hinzugekommenen Messungen wird das Ergebnis noch etwas kleiner, nämlich 1 zu 0.1278. Dieses Ergebnis ist erheblich kleiner als das Verhältnis, das man bisher anzunehmen pflegte, als man für die halbe Refraktion (gewöhnlich, wenn auch nicht ganz richtig, einfach Refraktion genannt) das 0.08-Fache der Krümmung des Erdbogens ansetzte. Diese Bestimmung, bei deren Beobachtungen fast durchgängig Heliotroplicht als Zielpunkt diente und bei der so viele große Entfernungen vorkommen (die größte beträgt $11\tfrac{1}{2}$ Meilen), scheint jedoch mehr Vertrauen zu verdienen als andere, bei denen schon das Anvisieren terrestrischer Zielpunkte keine so große Genauigkeit bei den beobachteten Zenitdistanzen erlaubte und die Entfernungen außerdem gewöhnlich viel kleiner waren. Bei kleinen Entfernungen sind die Anomalien auch nach diesen Erfahrungen immer am größten, besonders am Vormittag, um die Mittagszeit und am frühen Nachmittag. Zu diesen Zeiten erfährt nach diesen Beobachtungen das dicht über den Erdboden streichende Licht an sonnigen Sommertagen in flachen Gegenden und bei kleineren Entfernungen gewöhnlich eine negative Refraktion. --- Beiträge zu geografischen Längenbestimmungen aus beobachteten Sonnenfinsternissen und Sternbedeckungen von Herrn Prof. \textsc{Wurm}. --- Beobachtungen von Planetenoppositionen, Verfinsterungen der Jupitermonde und Sternbedeckungen, die Herr \textsc{Sniadecky} an der Sternwarte in Wilna angestellt hat. --- Verschiedene Beobachtungen von Herrn \textsc{Rümker} in Paramatta. Die wichtigsten davon, nämlich die Beobachtungen des \textsc{Enke}'schen Kometen im Juni 1822, wurden bereits in Stück 26 dieser Zeitschrift, Jahrgang 1823 [S. 442 dieses Bandes], mitgeteilt. --- Der hierauf folgende interessante Aufsatz von Dr. \textsc{Olbers} berührt die erhabenste Seite der Astronomie, nämlich die Frage, ob das Weltall Grenzen hat oder nicht. Im Bewusstsein unseres Unvermögens, den Schleier zu heben, der dieses Geheimnis verhüllt, beschäftigen wir uns doch gern mit der Klärung solcher Zweifel, die die Annahme, das Weltall sei unendlich, mit unseren sonstigen Einsichten in Widerspruch zu bringen scheinen könnten. Zu einem solchen Zweifel, wenigstens an der unendlichen Anzahl von Sonnen, führt die Überlegung, dass unter dieser Annahme \emph{jede} vom Auge ausgehende Richtung irgendwo in geringerer oder größerer Entfernung auf einen Sonnenkörper treffen und daher das ganze Himmelsgewölbe in nirgends unterbrochenem Sonnenglanz erscheinen würde. Dieser Schluss ist an sich vollkommen richtig. Ein gleichförmiger Glanz wäre jedoch nur möglich, wenn der Weltraum absolut durchsichtig wäre und das Licht daher, wenn es ihn durchdringt, gar nichts an Intensität verlöre. Diese Voraussetzung ist aber völlig unbewiesen und tatsächlich schon an sich höchst unwahrscheinlich. Dr. \textsc{Olbers} zeigt nun, dass eine \SourceLeaf{637} verhältnismäßig nur sehr geringe Undurchsichtigkeit angenommen zu werden braucht, um den scheinbaren Widerspruch der Erfahrung mit einer unendlichen Zahl von Sonnen aufzuheben. Ist die Schwächung des Lichts auf seinem Weg vom Sirius zu uns nur $\frac{1}{800}$, so würde eine hohle Kugel mit einem Radius von 30000 Siriusweiten, in deren Mitte wir uns befänden und deren ganze innere Fläche Sonnenglanz ausströmte, uns nur $\frac{1}{700000}$ so hell erscheinen wie der Himmelsgrund in einer klaren Vollmondnacht, also gewiss gar keinen merklichen Eindruck auf unsere Augen machen. --- Fixsternverzeichnis von Herrn \textsc{Pond}, mitgeteilt von Prof. \textsc{Tralles}. Der Sinn dieses Verzeichnisses und der daraus gezogenen Folgerung ist uns nicht ganz verständlich. --- Weitere Nachrichten über den \textsc{Pons}'schen (\textsc{Enke}'schen) Kometen von Prof. \textsc{Enke}. Die neuen Beobachtungen dieses Himmelskörpers im Jahr 1822 bestätigen das schon aus den früheren folgende bemerkenswerte Phänomen, dass neben den durch die Planetenstörungen erklärbaren Veränderungen der Umlaufzeit noch eine fortschreitende Abnahme derselben stattfindet, die von einem Umlauf zum anderen jetzt etwas über drei Stunden beträgt. Wahrscheinlich haben wir hierin die erste sichere Erfahrung vom tatsächlichen Vorhandensein eines Widerstands, den die Kometen bei ihrer Bewegung im Weltraum erfahren. Hoffentlich kann dieser Himmelskörper schon bei seiner nächsten Annäherung an die Sonne im August 1825 wieder in Europa beobachtet werden. Zumindest wird nach der Bearbeitung des ebenso geschickten wie unermüdlichen Rechners keinerlei Ungewissheit darüber bestehen, wo er zu suchen sein wird. --- Astronomische Beobachtungen an der Prager Sternwarte von den Herren \textsc{David} und \textsc{Bittner}. Ähnliche Beobachtungen in Kremsmünster von Herrn \textsc{Derflinger} und in Prag von Herrn \textsc{Hallaschka}. --- Originalbeobachtungen des dritten Kometen von 1822 von Dr. \textsc{Olbers}. --- Einige mechanische Untersuchungen über die Entstehungen der Kometenschweife von Dr. \textsc{Lehmann}. Dass alle Versuche, diesen rätselhaften Gegenstand befriedigend zu erklären, bisher ohne erheblichen Erfolg geblieben sind, darf uns nicht wundern. Erfahrungen von solcher Bestimmtheit, dass daraus exakte Ergebnisse als Grundlage strenger Rechnungen gewonnen werden könnten, fehlen uns gänzlich. Über die Beschaffenheit der dabei wirkenden Kräfte wissen wir wenig Gewisses, und die strenge mathematische Behandlung der Wirkung hypothetisch angenommener Kräfte bereitet sehr große Schwierigkeiten. Dieser Versuch von Herrn \textsc{Lehmann}, die Kometenschweife bloß aus den bekannten Kräften und ohne Hinzuziehung von Repulsivkräften zu erklären, bleibt zwar (wie alle anderen Erklärungsversuche) noch viel zu sehr an der Oberfläche des Problems und hält sich noch viel zu sehr an inadäquate Vorstellungen im Allgemeinen, als dass sich so auch nur über die Möglichkeit einer Erklärung auf diesem Weg urteilen ließe. Indessen ist er nicht ohne einfallsreiche Ansichten, und man kann nicht leugnen, dass jenes auch für die bisherigen Versuche mit Repulsivkräften gilt, deren Notwendigkeit noch nicht als entschieden betrachtet werden darf. --- Astronomische Beobachtungen an der königlichen Sternwarte in Berlin im Jahr 1822. --- Ein stärker vergrößernder Okularaufsatz für achromatische Fernrohre, erfunden von Herrn W. \textsc{Kitchiner} in London. So viel man aus der sehr unvollständigen und unklaren Angabe schließen kann, besteht die Erfindung nur darin, die Okularlinsen gegeneinander beweglich anzuordnen, um damit veränderliche Vergrößerungen zu erhalten. --- Beobachtungen und Elemente des Kometen vom September 1822 von Herrn \textsc{Rümker} in Paramatta. Diese Beobachtungen reichen bis etwa Mitte November, also beträchtlich weiter als die europäischen, und verdienen es daher, mit diesen verbunden zu werden. Ferner eine Beobachtung des Merkurdurchgangs am 5. November 1822 von demselben. --- Geografische Ortsbestimmungen in der Altmark und anderen Grenzen von Herrn Musikdirektor \textsc{Stöpel} in Tangermünde. Dieses wertvolle Verzeichnis enthält 134 Ortsbestimmungen nach trigonometrischen Messungen, die sich an die Dreiecksseite von Magdeburg zum Hagelsberg in der großen \textsc{von Müffling}'schen Triangulierung anschließen. Wir hätten gewünscht, dass Herr \textsc{Stöpel} zugleich einiges über seine Hilfsmittel und seine \SourceLeaf{638} Beobachtungsart mitgeteilt und sich nicht bloß deswegen auf ein auswärts unbekanntes Provinzialblatt bezogen hätte. --- Beobachtung einer Sternbedeckung durch Mars, von Prof. \textsc{Tralles}, umso bemerkenswerter, je seltener solche Erscheinungen sind. Der Stern wurde nach seinem Austritt erst in einer Entfernung von mehr als einem Planetendurchmesser wieder sichtbar, anfangs sehr schwach, dann aber schnell an Licht zunehmend. Das Fernrohr hatte $3\tfrac{1}{2}$ Zoll Öffnung. Herr \textsc{Tralles} schließt daraus auf das Dasein einer bedeutenden Marsatmosphäre. --- Über die von Herrn G. R. \textsc{Pastorf} entdeckte Photosphäre der Planeten, von Herrn \textsc{Ritz}. Herr \textsc{Ritz} erklärt diese Erscheinung aus einer doppelten Reflexion zwischen den Objektivlinsen. Der Rezensent hat immer dieselbe Ansicht gehabt und sich durch Rechnung nach den wirklichen Dimensionen von \textsc{Frauenhofer}'schen Fernrohren überzeugt, dass diese durch die gedachte Reflexion ein solches Phänomen wie das beobachtete erzeugen müssen. Eine wesentliche Bedingung der Entstehung ist, dass die Krümmungen der beiden einander zugekehrten Flächen \emph{beinahe} gleich sein müssen, wie es bei der Konstruktion der \textsc{Frauenhofer}'schen Objektive wirklich der Fall ist. --- Bemerkungen über den vorigen Gegenstand von Justizkommissionsrat \textsc{Kunovsky}. Herr \textsc{Kunovsky} gibt der Erklärungsart des Herrn \textsc{Ritz}, dessen Aufsatz ihm vor dem Abdruck zur Einsicht mitgeteilt worden zu sein scheint, seinen Beifall. Der Rezensent gesteht indessen, dass es ihm vorgekommen ist, als habe Herr \textsc{Kunovsky} den Geist der Erklärung nicht richtig aufgefasst, da beinahe alles, was er als eine Bestätigung anführt, wenn man es als faktisch bewiesen ansehen dürfte, in Widerspruch damit stehen oder wenigstens einer ganz anderen besonderen Erklärung bedürfen würde. Herr \textsc{Kunovsky} behauptet, die Erscheinung mit \emph{allen} Fernrohren, die er auf die Probe gestellt hat, bemerkt zu haben, obgleich bekannt ist, dass die englischen Objektive ganz anders konstruiert sind als die \textsc{Frauenhofer}'schen. Er findet die Lichtsphäre bei verschiedenen Fernrohren nahezu proportional zur Öffnung des Objektivs, was im Geist der obigen Erklärung nur von \emph{Einem und demselben} Fernrohr gilt, insofern dem Objektiv bald eine größere, bald eine geringere Öffnung gelassen wird. Bei verschiedenen Fernrohren, deren Dimensionen ungleich, aber in \emph{allen} Stücken einander proportional wären, müsste die Lichtsphäre durchaus gleich groß erscheinen. Endlich hat Herr \textsc{Kunovsky} bei Anwendung von viel stärkeren Vergrößerungen die Lichtsphäre des Jupiter, wie es sein muss, viel blasser werden sehen, während die Lichtsphäre von $\alpha$ in der Leyer durch stärkere Vergrößerung fast gar nichts von ihrer Helligkeit verloren haben soll, was auch wiederum mit obiger Erklärung unverträglich sein würde. Sehr brauchbar ist übrigens die von Herrn \textsc{Kunovsky} angedeutete Prüfung der Reinheit des Glases der Objektive. --- Aus einem Schreiben von Prof. \textsc{Bürg} werden einige Resultate seiner neueren Untersuchungen über die Mondtheorie mitgeteilt. --- Fragmente zur Erklärung des Aratus, von Prof. \textsc{Schaubach}. --- Beobachtung der totalen Mondfinsternis am 26. Januar 1823 und einer Sternbedeckung von Herrn \textsc{Rümker} im Paramatta. --- Einige Beobachtungen auf der Dorpater Sternwarte. --- Ähnliche auf der Sternwarte in Krakau von Prof. \textsc{Leski}. --- Über die astronomische Strahlenbrechung von Prof. \textsc{Bessel}. --- Astronomische Nachrichten von Herrn Prediger \textsc{Luthmer} in Hannover. \textsc{Herschel}s Grabschrift. --- Neue Elemente der Junobahn von Prof. \textsc{Nicolai}. Dies sind die sehr schätzbaren Resultate einer mühsamen, auf Quadraturen gegründeten Berechnung der Störungen, welche Jupiter seit 1804 auf die Bewegung der Juno ausgeübt hat, und der Diskussion von 15 seit der Entdeckung beobachteten Oppositionen. Diese werden durch die Rechnung mit großer Genauigkeit dargestellt. Die noch übrig bleibenden kleinen Unterschiede ist Herr \textsc{Nicolai} geneigt, einer von ihm vermuteten Unzulänglichkeit des \textsc{Newton}'schen Gravitationsgesetzes zuzuschreiben, indem er es nämlich für möglich hält, dass bei gleicher Entfernung Jupiter eine andere Attraktion auf die Sonne als auf die Juno ausübt. Für absolut unmöglich kann man freilich eine solche, auch schon von anderen aufgestellte Hypothese nicht erklären, allein sehr unwahrscheinlich ist sie doch schon deswegen, weil es faktisch bewiesen ist, dass die Attraktion der Sonne auf alle Haupt- und Nebenplaneten \SourceLeaf{639} gleich groß ist, ebenso wie die Gravitation der verschiedensten Körper auf der Erdoberfläche gegen die Erde. Nach Ansicht des Rezensenten darf man daher zu einer der Analogie so sehr widersprechenden Voraussetzung erst dann seine Zuflucht nehmen, wenn dafür entscheidende Beweise vorhanden sind, was bis jetzt keineswegs der Fall ist. --- Beobachtete Sternbedeckungen auf der Wiener Sternwarte von Prof. \textsc{Littrow}. --- Verzeichnis von 795 Doppelsternen aus Herrn \textsc{Struve}s Astr. Beobachtungen, Band 3. --- Unter den kürzeren astronomischen Nachrichten und Bemerkungen findet sich auch noch manches Interessante, obwohl sich zu verschiedenen Artikeln noch Einwände vorbringen ließen. So sieht z. B. der Rezensent nicht recht ein, inwiefern der Umstand, dass bei dem \textsc{Halley}'schen Kometen die Unterschiede der beobachteten Umlaufzeiten \emph{untereinander} viel größer sind als bei dem \textsc{Enke}'schen Kometen, einen Beweis für die großen Fortschritte der Beobachtungskunst und des Kalküls geben soll, so unleugbar letztere an sich sind. Selbst dass der Unterschied zwischen der \emph{berechneten} Umlaufzeit und der beobachteten bei dem \textsc{Halley}'schen Kometen viel größer war als bei dem \textsc{Enke}'schen, beweist an sich noch wenig für die Fortschritte des Kalküls, da die Umstände bei beiden Kometen so sehr verschieden waren.
